\RequirePackage{pdfmanagement}
\SetKeys[document/metadata]{pdfversion=1.7, pdfstandard=A-2u, lang=en-GB}
\input{preamble/pdfA-compliance/glyphtounicode}
\documentclass[11pt,a4paper]{report}

\newcommand*{\PaperTitle}{Multistability in bimolecular \((2,4,2)\) mass-action networks: there are no networks beyond the ten of Banaji, Boros and Hofbauer}
\newcommand*{\PaperTitlePlain}{Multistability in bimolecular (2,4,2) mass-action networks: there are no networks beyond the ten of Banaji, Boros and Hofbauer}
\newcommand*{\PaperShortTitle}{Bistable bimolecular \((2,4,2)\) networks}
\newcommand*{\PaperSubtitle}{chemistry-open-problems \textperiodcentered{} 25 October 2026 \textperiodcentered{} preliminary, not yet verified}
\newcommand*{\PaperCredits}{Solver: Claude Opus 5.5 (effort xhigh) \textperiodcentered{} Verification: pending}
\newcommand*{\PaperKeywords}{mass-action kinetics, multistability, bistability, boundary equilibria, Poincar\'e index, reaction networks}
\input{chemop-preamble}
\addbibresource{references.bib}

\providecommand{\R}{\mathbb{R}}
\newcommand{\Q}{\mathcal{Q}}
\newcommand{\Net}{\mathcal{N}}
\newcommand{\relint}{\operatorname{relint}}
\newcommand{\cone}{\operatorname{cone}}
\newcommand{\indx}{\operatorname{ind}}
\newcommand{\Trop}{\operatorname{tr}}
\newcommand{\rx}[2]{\(#1\to#2\)}

\begin{document}
\makepapertitle

\begin{paperabstract}
Banaji, Boros and Hofbauer found ten bimolecular mass-action networks with two species, four reactions and rank two
that admit two asymptotically stable equilibria in the closed nonnegative quadrant, the origin and a positive
equilibrium. They asked whether any other such network, not dynamically equivalent to one of the ten, has this
property, and noted that the answer requires studying boundary equilibria. We show that there is none. We
enumerate all 27\,195 networks of the class (5\,716 classes up to dynamical equivalence and exchange of the two
species) and prove that for every network and every choice of rate constants, two distinct asymptotically stable
equilibria must be the origin and a positive equilibrium. The key step is an index argument: if the origin is
asymptotically stable relative to the quadrant, the Poincar\'e indices of the positive equilibria sum to at most
zero. Hence a second, saddle-type positive equilibrium is needed, which happens only for the thirty fold networks of
the authors' Table~1, and exactly ten of these have a stable origin. The proof is computer-assisted only through exact
rational computations (sign patterns of explicit polynomials for 838 four-source classes, and certificates with
integer rate constants for the ten networks).
\end{paperabstract}

\section*{The question}

Banaji, Boros and Hofbauer~\autocite{BBH2024} (BBH below) study planar, quadratic mass-action networks with few
reactions. In §3.2 they list (Table~1) the thirty bimolecular \((2,4,2)\) networks that admit a fold bifurcation,
and they mark ten of them with an orange bullet: for these, the origin and a positive equilibrium can be
asymptotically stable at the same time (their Lemma~27). The last paragraph of their Conclusions (§7; p.~22 of
arXiv:2406.13451v1, p.~21446 of the published version) reads:

\begin{openquote}
``We have identified 10 bimolecular (2, 4, 2) networks (those marked with an orange bullet in Table 1) that admit
multiple nonnegative asymptotically stable equilibria. It is an open problem whether there are any further
bimolecular (2, 4, 2) mass-action networks with the same property, and which are not dynamically equivalent to any of
the 10 we have found. Answering this question requires studying boundary equilibria.''
\end{openquote}

We answer the question negatively.

\begin{theorem}\label{thm:main}
Let \(\Net\) be a bimolecular mass-action network with two species, four reactions and rank two. The following are
equivalent.
\begin{enumerate}
\item For some rate constants \(\kappa\in\R^4_{>0}\), the mass-action system of \(\Net\) has two distinct equilibria
in \(\R^2_{\ge0}\) that are asymptotically stable relative to \(\R^2_{\ge0}\).
\item \(\Net\) is dynamically equivalent, possibly after exchanging the two species, to one of the ten networks
marked with an orange bullet in Table~1 of~\autocite{BBH2024} (listed in \zcref{tab:ten}).
\end{enumerate}
Moreover, for every such network and every \(\kappa\), at most two equilibria in \(\R^2_{\ge0}\) are asymptotically
stable, and if two are, they are the origin and a positive equilibrium.
\end{theorem}

The ten classes contain 50 networks in total (\zcref{rem:fifty}). The proof combines elementary facts on boundary
equilibria (§2), a bound on positive equilibria (§3), and an index estimate that captures the boundary (§4). §5
assembles the proof and treats the thirty networks of Table~1. §6 lists the computer checks.

\section{Setting}

\paragraph{Networks.} The species are \(X,Y\) with concentrations \(z=(x,y)\). A complex \(aX+bY\) is identified
with \((a,b)\in\mathbb Z^2_{\ge0}\). It is bimolecular if \(a+b\le2\), so there are six of them:
\(0,X,Y,2X,X+Y,2Y\). A reaction \(s\to p\) is an ordered pair of distinct complexes, with reaction vector \(v=p-s\).
There are \(30\) bimolecular reactions. Following BBH (§2.1), a bimolecular \((2,4,2)\) network is a set of four
bimolecular reactions whose reaction vectors span \(\R^2\). There are \(27\,195\) of them.

\paragraph{Mass action.} With rate constants \(\kappa\in\R^4_{>0}\), the system is
\begin{equation}\label{eq:ma}
\dot z=F(z):=\sum_{j=1}^4\kappa_j\,z^{s_j}v_j=\sum_{s}z^s\,w_s,\qquad w_s:=\sum_{j:\,s_j=s}\kappa_jv_j,
\end{equation}
where \(z^{(a,b)}=x^ay^b\) and the second sum runs over the distinct sources of~\(\Net\). Note two sign rules:
\(v_1\ge-s_1\) and \(v_2\ge-s_2\) (a reaction cannot consume more than its source contains), and
\(v_1+v_2\le 2-s_1-s_2\). Hence \(F_1(0,y)\ge0\) and \(F_2(x,0)\ge0\), so the closed quadrant
\(\Q:=\R^2_{\ge0}\) and the open quadrant are forward invariant.

\paragraph{Dynamical equivalence.} BBH (§2.2, following~\autocite{CraciunPantea2008}) call two networks
dynamically equivalent if they give rise to the same set of mass-action differential equations, perhaps after
permuting species and reactions. As \(\kappa\) ranges over \(\R^4_{>0}\), the coefficient \(w_s\) in~\eqref{eq:ma}
ranges over the strictly positive combinations of the reaction vectors with source \(s\), which form the relative
interior of the closed convex cone \(C_s\) they generate~\autocite[Cor.~6.6.2]{Rockafellar1970}, independently for
each source. Since distinct monomials are linearly independent functions, two networks on \(\{X,Y\}\) give the same
set of vector fields if and only if they have the same sources and the same cones \(C_s\). We use this as the
definition of the classes, and also identify networks that differ by the exchange \(X\leftrightarrow Y\). There are
\(5\,716\) classes. A network is \emph{dynamically nontrivial} if \(\Gamma\kappa=0\) for some \(\kappa>0\), where
\(\Gamma\) is the \(2\times4\) matrix of reaction vectors. Exactly \(2\,103\) classes are nontrivial, and \(838\) of
them have four distinct sources, as in BBH's Lemma~24.

\paragraph{Stability.} An equilibrium \(e\in\Q\) is \emph{asymptotically stable relative to \(\Q\)} (AS) if for every
\(\varepsilon>0\) there is \(\delta>0\) such that every solution starting in \(\Q\cap B_\delta(e)\) stays in
\(B_\varepsilon(e)\) and converges to \(e\). This is the notion used by BBH in the proof of their Lemma~27, where the
origin of three of the ten networks is AS although the vector field has a saddle-node there. For positive equilibria
it is ordinary asymptotic stability. An AS equilibrium is isolated among the equilibria in \(\Q\). We say that
\(\Net\) has property~(MS) if (1) of \zcref{thm:main} holds. (MS) depends only on the class of~\(\Net\).

\paragraph{Index.} For an isolated zero \(e\) of a continuous planar vector field, \(\indx(e)\) is the winding number
of the field along a small circle around~\(e\). If \(\det DF(e)\neq0\), then \(\indx(e)=\operatorname{sign}\det
DF(e)\). For a positively oriented Jordan curve on which \(F\neq0\), the winding number of \(F\) along the curve is
the sum of the indices of the (finitely many) zeros inside.

\section{Boundary equilibria}

Write \(F_2(x,0)=r(x)\) and \(F_1(x,0)=p(x)=p_0+p_1x+p_2x^2\). Here \(p_0,p_1,p_2\) are the first components of
\(w_0,w_X,w_{2X}\) (zero if the source is absent), so \(p_0\ge0\ge p_2\) by the sign rules.

\begin{lemma}\label{lem:axis}
Fix \(\Net\) and \(\kappa\).
\begin{enumerate}
\item The origin is an equilibrium if and only if \(\Net\) has no reaction with source~\(0\).
\item If \((x^*,0)\) with \(x^*>0\) is an equilibrium, then no reaction with source \(0\), \(X\) or \(2X\) has \(Y\) in
its product. Hence \(r\equiv0\), and \(F_2=y\,G(x,y)\) with \(G=a+bx+ey\), where \(a,b,e\) are the second
components of \(w_Y,w_{X+Y},w_{2Y}\). Moreover \(e<0\) if \(2Y\) is a source, and \(e=0\) otherwise.
\item If moreover \((x^*,0)\) is AS, then \(p\) changes sign from \(+\) to \(-\) at \(x^*\), and \(G(x^*,0)\le0\).
\item If the origin is AS, then it is the only equilibrium on the boundary of \(\Q\).
\end{enumerate}
The statements with \(X\) and \(Y\) exchanged also hold.
\end{lemma}

\begin{proof}
(1) \(F(0)=\sum\kappa_jv_j\) over the reactions with source \(0\), and each such \(v_j=p_j\) is nonzero and
nonnegative. (2) \(r(x^*)\) is a sum of nonnegative terms \(\kappa_j(x^*)^{s_{j,1}}p_{j,2}\) over the reactions with
source \(0,X,2X\), so each \(p_{j,2}\) vanishes. The other reactions have sources \(Y,X+Y,2Y\) and contribute
\(y,xy,y^2\) monomials. A reaction with source \(2Y\) has \(v_2\in\{-1,-2\}\). (3) The \(x\)-axis is invariant and
contained in \(\Q\), so \(x^*\) is AS for \(\dot x=p(x)\). If \(p\equiv0\), the axis consists of equilibria.
Otherwise \(p\) has degree at most two, and asymptotic stability forces a sign change from \(+\) to \(-\). If
\(G(x^*,0)>0\), then \(G>0\) on a neighbourhood \(U\) of \((x^*,0)\). A solution from \((x^*,y_0)\), \(y_0>0\) small,
has increasing \(y\) while it stays in \(U\), so it cannot both stay in \(U\) and converge to \((x^*,0)\).
(4) Let \((x^*,0)\), \(x^*>0\), be an equilibrium. By (1) and (2), \(p_0=0\), the \(x\)-axis is invariant and
\(p=x(p_1+p_2x)\). If \(p\equiv0\), the origin is not isolated. Otherwise \(x^*=-p_1/p_2\) with \(p_2\le0\) forces
\(p_1>0>p_2\), so \(\dot x=p(x)>0\) on \((0,x^*)\): points \((x,0)\) with small \(x>0\) move away from the origin,
which is then not AS. The \(y\)-axis is symmetric.
\end{proof}

\begin{proposition}\label{prop:bdry}
For every network and every \(\kappa\), at most one equilibrium on the boundary of \(\Q\) is AS.
\end{proposition}

\begin{proof}
Up to the exchange of species, two boundary equilibria are of one of three kinds.

\emph{Two points of the open \(x\)-axis.} By \zcref{lem:axis}(3), \(p\not\equiv0\) changes sign from \(+\) to \(-\)
at both, hence also from \(-\) to \(+\) in between. Then \(p\) has three zeros, which is impossible for
\(p\not\equiv0\) of degree at most two.

\emph{The origin and \((x^*,0)\).} Excluded by \zcref{lem:axis}(4).

\emph{\((x^*,0)\) and \((0,y^*)\).} Both axes are invariant. A reaction with source \(0\) would have \(X\) or \(Y\) in
its product, contradicting \zcref{lem:axis}(2) or its mirror image. So \(p_0=0\), \(p=x(p_1+p_2x)\), and the sign
change at \(x^*\) forces \(p_1>0>p_2\).
Thus \(\Net\) contains \rx{X}{2X} (the only reaction with source \(X\) along the \(x\)-axis that increases \(x\)) and a
reaction with source \(2X\). By symmetry it contains \rx{Y}{2Y} and a reaction with source \(2Y\). These are the
four reactions, so \(F=(x(p_1+p_2x),\,y(q_1+q_2y))\) with \(q_1>0\). Then \(G(x^*,0)=q_1>0\), contradicting
\zcref{lem:axis}(3).
\end{proof}

\begin{proposition}\label{prop:axispos}
For no network and no \(\kappa\) are an equilibrium \((x^*,0)\) with \(x^*>0\) and a positive equilibrium both AS.
The same holds for \((0,y^*)\).
\end{proposition}

\begin{proof}
Suppose both are AS. By \zcref{lem:axis}, the reactions with sources \(0,X,2X\) lie along the \(x\)-axis, and
\(p\) changes sign at \(x^*\). Then at least two of \(p_0,p_1,p_2\) are nonzero: a single monomial has no positive
zero. So \(\Net\) has reactions with at least two distinct sources among \(0,X,2X\). These reactions do not change
\(y\). Since the rank is two, at least one reaction changes \(y\), and its source is \(Y\), \(X+Y\) or \(2Y\). Hence
there are one or two reactions with sources in \(\{Y,X+Y,2Y\}\), and the positive equilibria lie on \(\{G=0\}\).

Because \(p_0\ge0\ge p_2\), \(p\) has exactly one positive zero, so \(p>0\) on \((0,x^*)\) and \(p<0\) on
\((x^*,\infty)\).

\emph{All reactions with sources in \(\{Y,X+Y,2Y\}\) share one source.} Then \(G\) is a single monomial times a
coefficient. If the coefficient is nonzero, \(G\) has no zero in \(\R^2_{>0}\). If it is zero, \(\dot y\equiv0\):
every horizontal line is invariant, so no positive equilibrium attracts the points above or below it.

\emph{Sources \(Y\) and \(X+Y\), one reaction each.} Then \(G=a+bx\), and a positive zero needs \(ab<0\), at
\(\bar x=-a/b\). Write \(F_1=p(x)+y(\alpha+\beta x)\), where \(\alpha,\beta\) are the first components of
\(w_Y,w_{X+Y}\). If \(\alpha+\beta\bar x=0\), the equilibria on \(x=\bar x\) form a vertical segment or do not exist.
Otherwise the positive equilibrium is \((\bar x,\bar y)\) with \(\bar y=-p(\bar x)/(\alpha+\beta\bar x)>0\), and
\[
DF(\bar x,\bar y)=\begin{pmatrix}p'(\bar x)+\beta\bar y & \alpha+\beta\bar x\\ b\bar y & 0\end{pmatrix},\qquad
\det DF=-b\bar y(\alpha+\beta\bar x)=b\,p(\bar x).
\]
Asymptotic stability requires \(\det DF\ge0\). Equality would mean \(p(\bar x)=0\), hence \(\bar y=0\). So
\(b\,p(\bar x)>0\). On the other hand, \zcref{lem:axis}(3) gives \(G(x^*,0)=b(x^*-\bar x)\le0\). If \(b>0\), then
\(x^*\le\bar x\) and \(p(\bar x)\le0\). If \(b<0\), then \(x^*\ge\bar x\) and \(p(\bar x)\ge0\). Both contradict
\(b\,p(\bar x)>0\).

\emph{Sources \(Y\) and \(2Y\).} Here \(G=a+ey\) with \(e<0\). A positive zero needs \(a>0\), and then
\(G(x^*,0)=a>0\).

\emph{Sources \(X+Y\) and \(2Y\).} Here \(G=bx+ey\) with \(e<0\). A positive zero needs \(b>0\), and then
\(G(x^*,0)=bx^*>0\).

In the last two cases \zcref{lem:axis}(3) is violated.
\end{proof}

\section{Positive equilibria}

\begin{lemma}\label{lem:few}
Fix \(\Net\) and \(\kappa\).
\begin{enumerate}
\item If \(\Net\) is dynamically trivial, there is no positive equilibrium.
\item If \(\Net\) has at most two distinct sources, no positive equilibrium is isolated.
\item If \(\Net\) has three distinct sources, the set of positive equilibria is empty, a single point, or has no
isolated point.
\end{enumerate}
\end{lemma}

\begin{proof}
(1) \(F(z)=\Gamma(\kappa\circ z^{A})\) with \(\kappa\circ z^A>0\). (2) With one source, \(F=z^sw_s\) vanishes nowhere
or everywhere. With two sources, \(F=z^{s'}(z^{s-s'}w_s+w_{s'})\) vanishes exactly where \(z^{s-s'}w_s=-w_{s'}\).
This set is empty, the whole quadrant, or a level curve of the monomial \(z^{s-s'}\).

(3) Let \(a_1,a_2,a_3\) be the sources and \(W=[w_1\,w_2\,w_3]\). If \(\operatorname{rank}W=2\), then \(F(z)=0\) iff
\((z^{a_i})_i\) is a multiple of the kernel vector \(k\) of \(W\). This requires \(k>0\) (up to sign). Then
\(z^{a_i-a_1}=k_i/k_1\) for \(i=2,3\). If \(a_2-a_1\) and \(a_3-a_1\) are independent, this log-linear system has
exactly one solution. Otherwise its solution set is empty or a level curve. If \(\operatorname{rank}W\le1\), then
\(F=\phi(z)\,u\) with \(\phi=\sum c_iz^{a_i}\) and a fixed \(u\). In logarithmic coordinates \(\zeta=\ln z\), write
\(\phi=e^{\langle a_1,\zeta\rangle}\psi(\zeta)\) with \(\psi=c_1+c_2e^{\langle a_2-a_1,\zeta\rangle}+c_3e^{\langle
a_3-a_1,\zeta\rangle}\). If \(a_2-a_1,a_3-a_1\) are independent, \(\nabla\psi\neq0\) wherever \(\psi=0\), unless
\(c_2=c_3=0\). So the zero set of \(\psi\) is a curve, empty, or everything. If they are parallel, \(\psi\) depends on
one linear form, and its zero set is a union of lines.
\end{proof}

For four distinct sources, BBH's parametrisation (their §2.4--2.6 and Lemma~11) describes the positive
equilibria exactly. We recall it in the form we need. Let \(\Net\) be nontrivial with four distinct sources
\(a_1,\dots,a_4\) (one reaction each), \(A\) the \(4\times2\) matrix with rows \(a_i\), and \(\Gamma=[v_1\cdots
v_4]\). The cone \(\ker\Gamma\cap\R^4_{\ge0}\) is two-dimensional and pointed, with extreme rays
\(u,v\in\mathbb Z^4_{\ge0}\). Put \(h(t)=(1-t)u+tv\) for \(t\in I:=(0,1)\). The four sources are never collinear (no
line contains four of the six complexes), so \([A\,|\,\mathbf 1]\) has rank three, with left kernel spanned by an
integer vector \(c\). Define
\[
\Phi(t)=\sum_ic_i\ln h_i(t),\qquad N(t)=\Phi'(t)\prod_ih_i(t),\qquad Q(t)=\det\bigl(\Gamma\operatorname{diag}(h(t))A\bigr).
\]
\(N\) and \(Q\) are polynomials with integer coefficients that do not depend on \(\kappa\).

\begin{proposition}\label{prop:four}
Let \(\Net\) be as above.
\begin{enumerate}
\item For each \(\kappa\), the positive equilibria correspond bijectively to the roots \(t\in I\) of
\(\Phi(t)=c\cdot\ln\kappa\). At the equilibrium attached to \(t\), \(\det DF=\lambda^2Q(t)/(xy)\) for some
\(\lambda>0\).
\item (Exact computation, \texttt{positive\_equilibria.py}.) For 9 of the 838 classes, \(N\equiv Q\equiv0\). For the
other 829, \(QN\) has no sign change on \(I\). Moreover, either \(Q\) has no zero in \(I\), or it has exactly one zero
in \(I\), which is simple. The second case occurs exactly for the 30 networks of BBH's Table~1.
\end{enumerate}
\end{proposition}

\begin{proof}
(1) \(z>0\) is an equilibrium iff \(\kappa\circ z^A\in\ker\Gamma\cap\R^4_{>0}=\{\lambda h(t):\lambda>0,t\in I\}\),
i.e. iff \(\ln\kappa+A\ln z=\ln\lambda\,\mathbf 1+\ln h(t)\) for some \(\lambda,t\). Pairing with \(c\) gives
\(c\cdot\ln\kappa=\Phi(t)\). Conversely, if \(\Phi(t)=c\cdot\ln\kappa\), then \(\ln h(t)-\ln\kappa\) is orthogonal to
\(c\), hence lies in the image of \([A\,|\,\mathbf1]\). Since this matrix is injective, \((\ln z,\ln\lambda)\) is
unique. Distinct \(t\) give distinct \(z\), because \(h(t)\) and \(h(t')\) are not proportional. Finally
\(DF(z)=\Gamma\operatorname{diag}(\kappa\circ z^A)A\operatorname{diag}(1/z)=\lambda\Gamma\operatorname{diag}(h(t))A
\operatorname{diag}(1/z)\).
\end{proof}

\begin{corollary}\label{cor:pos}
For every network and every \(\kappa\):
\begin{enumerate}
\item at most one positive equilibrium is AS;
\item if there are at least two positive equilibria and one of them is isolated, then \(\Net\) has four distinct
sources and belongs to one of the 30 classes of Table~1 of~\autocite{BBH2024}. The set of positive equilibria is then
finite.
\end{enumerate}
\end{corollary}

\begin{proof}
An AS positive equilibrium is isolated. By \zcref{lem:few}, two isolated positive equilibria require four distinct
sources and nontriviality. In the 9 classes with \(N\equiv0\), \(\Phi\) is constant, so there is no isolated positive
equilibrium. Assume \(N\not\equiv0\). Then \(\Phi(t)=c\cdot\ln\kappa\) is equivalent to a nonzero polynomial
equation, so the positive equilibria are finitely many. Since \(QN\) has no sign change on \(I\), \(\Phi'\) changes
sign exactly where \(Q\) does. If \(Q\) has no zero in \(I\), then \(\Phi\) is strictly monotone and there is at most
one positive equilibrium. Otherwise \(Q\) has a single, simple zero \(\beta\in I\), and \(\Phi\) is strictly monotone
on \((0,\beta]\) and on \([\beta,1)\), in opposite directions. So there are at most two positive equilibria. If there
are two, at \(t_1<\beta<t_2\), then \(Q(t_1)\) and \(Q(t_2)\) are nonzero with opposite signs. The one with
\(\det DF<0\) is a saddle, which is not AS.
\end{proof}

This recovers, in exact arithmetic, BBH's Theorem~25 and Remark~26(d), which they obtained with a Mathematica
notebook~\autocite{BorosRepo}: outside Table~1, a bimolecular \((2,4,2)\) network has at most one positive
equilibrium for each \(\kappa\).

\section{The index of a stable origin}

\begin{lemma}\label{lem:kras}
Let \(F\) be a \(C^1\) vector field on an open set of \(\R^2\).
\begin{enumerate}
\item If \(e\) is an asymptotically stable equilibrium, then \(\indx(e)=1\).
\item Suppose \(0\) is an equilibrium, \(\Q\) is forward invariant near \(0\), and \(0\) is AS relative to \(\Q\).
For small \(\varepsilon>0\), let \(\gamma\) be the arc \(\{z\in\Q:|z|=\varepsilon\}\), run clockwise from
\(A=(0,\varepsilon)\) to \(B=(\varepsilon,0)\). Then \(F\neq0\) on \(\gamma\). Let \(\theta_A\in[-\pi/2,\pi/2]\) and
\(\theta_B\in[0,\pi]\) be the angles of \(F(A)\) and \(F(B)\); these ranges are possible because \(F_1(A)\ge0\) and
\(F_2(B)\ge0\). Then the change of a continuous angle of \(F\) along \(\gamma\) equals \(\theta_B-\theta_A-2\pi\).
\end{enumerate}
\end{lemma}

\begin{proof}
(1) is classical~\autocite{KrasnoselskiiZabreiko1984}. It follows from the argument below with \(\gamma\) replaced
by a full circle around \(e\): there \(H(\cdot,T)\) points strictly inwards, hence is homotopic, without zeros on
the circle, to \(z\mapsto e-z\), of index \(1\). We prove (2). Let \(\varphi_t\)
be the flow. Choose \(\varepsilon\) so small that every solution from \(\Q\cap\bar B_\varepsilon\) converges to \(0\).
Then \(\varphi_\tau(z)\neq z\) for \(z\in\gamma\) and \(\tau>0\), since such a \(z\) would be an equilibrium or a
periodic point different from \(0\). By compactness of \(\gamma\) there is \(T>0\) with \(|\varphi_T(z)|<\varepsilon/2\) for
\(z\in\gamma\). Put \(H(z,\tau)=(\varphi_\tau(z)-z)/\tau\) for \(0<\tau\le T\) and \(H(z,0)=F(z)\). This is continuous
and nonzero on \(\gamma\times[0,T]\). Because \(\Q\) is forward invariant, \(H_1(A,\tau)\ge0\) and
\(H_2(B,\tau)\ge0\). So the quantity \(\omega(\tau)=\Delta_\gamma(\tau)-(\theta_B(\tau)-\theta_A(\tau))\) is defined,
with the end angles taken in the closed half-planes as in the statement. It is continuous, and it lies in
\(2\pi\mathbb Z\), so it is constant. At \(\tau=T\) we have \(H(z,T)\cdot z<(\varepsilon^2/2-\varepsilon^2)/T<0\) on
\(\gamma\). If \(z\) has polar angle \(\psi\in[0,\pi/2]\), the angle of \(H(z,T)\) can thus be chosen continuously in
\((\psi+\pi/2,\psi+3\pi/2)\). At \(A\) (\(\psi=\pi/2\)) this interval meets \(\theta_A+2\pi\mathbb Z\) only in
\(\theta_A+2\pi\). At \(B\) (\(\psi=0\)) it meets \(\theta_B+2\pi\mathbb Z\) only in \(\theta_B\). Hence
\(\omega(T)=-2\pi\).
\end{proof}

\begin{lemma}\label{lem:sweep}
Let \(c:[\sigma_0,\sigma_1]\to\R^2\setminus\{0\}\) have polynomial components of degree at most two, and let
\(\theta\) be a continuous angle of \(c\). Then \(\max\theta-\min\theta\le2\pi\).
\end{lemma}

\begin{proof}
If \(\max\theta-\min\theta>2\pi\), pick \(\vartheta\) with \(\vartheta,\vartheta+\pi,\vartheta+2\pi\) in
\((\min\theta,\max\theta)\). By the intermediate value theorem, \(c\) meets the line \(\R(\cos\vartheta,\sin\vartheta)\)
at three distinct parameters. So the polynomial \(\det(c(\sigma),(\cos\vartheta,\sin\vartheta))\), of degree at most
two, has three zeros and vanishes identically. Then \(c\) stays on that line and, avoiding \(0\), on one half of it.
So \(\theta\) is constant, a contradiction.
\end{proof}

\begin{theorem}\label{thm:index}
Let \(F\) be a polynomial vector field of degree at most two with \(F_1(0,y)\ge0\) and \(F_2(x,0)\ge0\) for
\(x,y\ge0\). Suppose the origin is an equilibrium that is AS relative to \(\Q\), that \(F\) has no other zero on the
boundary of \(\Q\), and that \(F\) has finitely many zeros in \(\R^2_{>0}\). Then the sum of their indices is at most
\(0\).
\end{theorem}

\begin{proof}
Take \(\varepsilon\) as in \zcref{lem:kras}(2), and \(M>0\) larger than \(x+y\) at every positive zero. Let \(D\) be
the region \(\{z\in\Q:x+y\le M,|z|\ge\varepsilon\}\), whose boundary, run counterclockwise, consists of four pieces:
the segment from \(B=(\varepsilon,0)\) to \((M,0)\), the segment \(S\) from \((M,0)\) to \((0,M)\), the segment from
\((0,M)\) to \(A=(0,\varepsilon)\), and the arc \(\gamma\) from \(A\) to \(B\). Then \(F\neq0\) on \(\partial D\), and
the sum \(W\) of the indices of the positive zeros satisfies \(2\pi W=\) total change of the angle of \(F\) along
\(\partial D\).

On the \(x\)-axis, \(F\) lies in the closed upper half-plane, where the angle can be taken in \([0,\pi]\). So the
change along the first segment is \(\theta(M,0)-\theta_B\), with \(\theta(M,0)\in[0,\pi]\). Similarly the change
along the \(y\)-axis segment is \(\theta_A-\theta(0,M)\), with \(\theta(0,M)\in[-\pi/2,\pi/2]\). Along \(\gamma\) the
change is \(\theta_B-\theta_A-2\pi\) by \zcref{lem:kras}(2). Along \(S\), parametrised by \(\sigma\mapsto
F(M-\sigma,\sigma)\), the components are polynomials of degree at most two. By \zcref{lem:sweep}, the change
\(\Delta_S\) satisfies \(|\Delta_S|\le2\pi\), and \(\Delta_S=\delta+2\pi n\) with \(n\in\mathbb Z\) and
\(\delta:=\theta(0,M)-\theta(M,0)\in[-3\pi/2,\pi/2]\). Summing,
\[
2\pi W=\Delta_S-\delta-2\pi=2\pi(n-1),\qquad 2\pi n=\Delta_S-\delta\le2\pi+\tfrac{3\pi}{2},
\]
so \(n\le1\) and \(W\le0\).
\end{proof}

\section{Proof of the main theorem}

\begin{proof}[of \zcref{thm:main}]
\emph{(1)\(\Rightarrow\)(2).} Let \(E_1\neq E_2\) be AS equilibria in \(\Q\) for some \(\kappa\). By
\zcref{prop:bdry}, one of them, say \(E_2\), is positive. By \zcref{cor:pos}(1), \(E_1\) is not positive. By
\zcref{prop:axispos}, \(E_1\) is not on an open half-axis. So \(E_1\) is the origin. By \zcref{lem:axis}(1),
\(\Net\) has no reaction with source \(0\), and by \zcref{lem:axis}(4) there is no other boundary equilibrium. The set of
positive equilibria is finite: it contains the isolated point \(E_2\), so by \zcref{lem:few} \(\Net\) has three or
four sources, and \zcref{lem:few}(3) or the proof of \zcref{cor:pos} applies. \zcref{thm:index} then gives
\(\sum\indx\le0\) over the positive equilibria. Since \(\indx(E_2)=1\) by \zcref{lem:kras}(1), there is a second
positive equilibrium. By \zcref{cor:pos}(2), \(\Net\) is one of the 30 networks of Table~1 of~\autocite{BBH2024}, and
it has no inflow reaction. \zcref{tab:thirty} sorts these 30 networks. Fifteen contain an inflow reaction \(0\to C\).
Five have no inflow, but their origin is never AS:
\begin{itemize}
\item \(X\to2X\), \(X+Y\to0\), \(Y\to2Y\), with \(2Y\to X\) or \(2Y\to2X\). Here the \(x\)-axis is invariant and
\(\dot x=\kappa_1x\) on it, so points of the axis leave every neighbourhood of~\(0\).
\item \(Y\to2X\), \(X+Y\to2Y\), \(2X\to0\), with \(X\to Y\), \(X\to X+Y\) or \(X\to2Y\). Here \(DF(0)\) has positive
off-diagonal entries and \(\det DF(0)=-\kappa\kappa'\), \(-2\kappa\kappa'\) or \(-3\kappa\kappa'\) (rate constants of
\(Y\to2X\) and of the reaction with source \(X\)). So \(0\) is a hyperbolic saddle, and its unstable eigenvector is
the Perron vector of a nonnegative irreducible matrix, hence positive. The unstable manifold therefore enters
\(\R^2_{>0}\)~\autocite[§2.7]{Perko2001}, and its points leave every small neighbourhood of~\(0\).
\end{itemize}
The remaining ten are the networks of (2).

\emph{(2)\(\Rightarrow\)(1).} \zcref{tab:ten} gives, for each of the ten networks, integer rate constants for which
\((1,1)\) is an equilibrium with \(\det DF(1,1)>0>\Trop DF(1,1)\), and the origin is AS. For seven of them
\(\det DF(0)>0>\Trop DF(0)\). For the three with \(2X\to0\), \(DF(0)\) is singular, and we use \(V=x+3y\). Along
solutions, \(\dot V=\ell y+Dx^2+Exy+Fy^2\) with \(\ell<0\) and \(D<0\), with no constant and no \(x\) term (exact
values in \texttt{certify\_ten.log}). On \(\Q\), \(\dot V\le y(\ell+|E|x+\max(F,0)y)+Dx^2<0\) whenever \(z\neq0\) and
\(|E|x+\max(F,0)y<-\ell\). The compact triangles \(\{z\in\Q:V(z)\le\rho\}\), \(\rho\) small, are therefore forward
invariant, shrink to \(0\), and every solution in them converges to \(0\) (LaSalle). Since \(\Q\) is forward
invariant and \(V\) decreases, the origin is AS relative to~\(\Q\).

\emph{Last statement.} By \zcref{prop:bdry} and \zcref{cor:pos}(1), at most one boundary equilibrium and at most
one positive equilibrium are AS. The pair is (origin, positive) by the argument above.
\end{proof}

\begin{table}[ht]
\centering
\caption{The ten networks of BBH's Table~1 with an orange bullet, with the integer rate constants (in the order
listed) certified by \texttt{scripts/certify\_ten.py}. For these constants the origin and \((1,1)\) are both AS, and
there is exactly one other positive equilibrium, a saddle. Type H: \(DF(0)\) is a sink. Type L: \(DF(0)\) is singular,
and stability follows from the Lyapunov function \(x+3y\).}\label{tab:ten}
\small
\begin{tabular}{lll}
\hline
Network & \(\kappa\) & origin \\
\hline
\(Y\to2X\), \(X+Y\to2Y\), \(X\to0\), \(2Y\to0\) & 81, 115, 47, 17 & H \\
\(Y\to2X\), \(X+Y\to2Y\), \(X\to0\), \(2Y\to X\) & 8, 14, 5, 3 & H \\
\(Y\to2X\), \(X+Y\to2Y\), \(2X\to0\), \(2Y\to0\) & 64, 86, 21, 11 & L \\
\(Y\to2X\), \(X+Y\to2Y\), \(2X\to0\), \(2Y\to X\) & 65, 191, 1, 63 & L \\
\(Y\to2X\), \(X+Y\to2Y\), \(2X\to0\), \(2Y\to2X\) & 2, 128, 1, 63 & L \\
\(Y\to2X\), \(2X\to2Y\), \(X\to0\), \(2Y\to0\) & 64, 41, 46, 9 & H \\
\(Y\to2X\), \(2X\to2Y\), \(X\to0\), \(2Y\to X\) & 64, 49, 47, 17 & H \\
\(Y\to2X\), \(2X\to2Y\), \(X\to0\), \(X+Y\to0\) & 83, 51, 45, 19 & H \\
\(Y\to2X\), \(2X\to2Y\), \(X\to0\), \(X+Y\to X\) & 65, 64, 2, 63 & H \\
\(Y\to2X\), \(2X\to2Y\), \(X\to0\), \(X+Y\to Y\) & 64, 32, 1, 63 & H \\
\hline
\end{tabular}
\end{table}

\begin{table}[ht]
\centering
\caption{The 30 networks of BBH's Table~1 (fold bifurcation) and the status of their origin. Each row lists the
two common reactions of a group, then the possible third reaction and the possible fourth reactions.}\label{tab:thirty}
\small
\begin{tabular}{p{3.2cm}p{5.6cm}p{3.0cm}}
\hline
Common reactions & Third reaction; fourth reaction & Origin \\
\hline
\(X\to2X\), \(X+Y\to0\) & \(0\to Y\); any of the 8 in BBH's group~1 & not an equilibrium (8) \\
\(X\to2X\), \(X+Y\to0\) & \(Y\to2Y\); \(2Y\to X\) or \(2Y\to2X\) & unstable, invariant axis (2) \\
\(Y\to2X\), \(X+Y\to2Y\) & \(X\to0\); \(0\to Y\) or \(0\to X+Y\) & not an equilibrium (2) \\
\(Y\to2X\), \(X+Y\to2Y\) & \(X\to0\); \(2Y\to0\) or \(2Y\to X\) & AS, hyperbolic (2) \\
\(Y\to2X\), \(X+Y\to2Y\) & \(2X\to0\); \(0\to Y\) or \(0\to X+Y\) & not an equilibrium (2) \\
\(Y\to2X\), \(X+Y\to2Y\) & \(2X\to0\); \(X\to Y\), \(X\to X+Y\) or \(X\to2Y\) & saddle (3) \\
\(Y\to2X\), \(X+Y\to2Y\) & \(2X\to0\); \(2Y\to0\), \(2Y\to X\) or \(2Y\to2X\) & AS, singular \(DF(0)\) (3) \\
\(Y\to2X\), \(2X\to2Y\) & \(X\to0\); \(0\to X\), \(0\to Y\) or \(0\to X+Y\) & not an equilibrium (3) \\
\(Y\to2X\), \(2X\to2Y\) & \(X\to0\); \(2Y\to0\), \(2Y\to X\), \(X+Y\to0\), \(X+Y\to X\) or \(X+Y\to Y\) & AS, hyperbolic (5) \\
\hline
\end{tabular}
\end{table}

\begin{remark}\label{rem:fifty}
Counting networks rather than classes, exactly 50 of the 27\,195 bimolecular \((2,4,2)\) networks have property
(MS). For example, Wilhelm's network \(Y\to2X\), \(2X\to X+Y\), \(X+Y\to Y\), \(X\to0\)~\autocite{Wilhelm2009} lies in
the class of \(Y\to2X\), \(2X\to2Y\), \(X\to0\), \(X+Y\to Y\). The members of each class are listed in
\texttt{certificates/classes.csv}.
\end{remark}

\begin{remark}
If asymptotic stability of a boundary equilibrium were required in a full neighbourhood in \(\R^2\), for the
polynomial vector field extended to the plane, only seven classes would remain. For the three networks with \(2X\to0\),
the origin is then a saddle-node and is not stable in the plane. Our negative result holds a fortiori for that
stronger notion.
\end{remark}

\begin{remark}
The theorem shows that, in this class, bistability always comes with two positive equilibria, i.e. with a fold
bifurcation. The index estimate (\zcref{thm:index}) is the study of boundary equilibria that BBH anticipated. It
uses only that \(F\) is quadratic with \(\Q\) forward invariant, so it holds for every planar mass-action system with
reactant complexes of molecularity at most two (any number of reactions, any product molecularity). In such a
system, if the origin is AS relative to \(\Q\) and the only boundary equilibrium, there cannot be exactly one positive
equilibrium that is asymptotically stable.
\end{remark}

\section{Computer verification}

All scripts are in \texttt{scripts/} and use exact rational arithmetic (sympy~1.14), except the sanity check.
Each prints PASS/FAIL per claim and exits with a nonzero status on failure. Logs are in \texttt{certificates/}.

\begin{itemize}
\item \texttt{networks.py}: enumeration and canonical forms (cones per source, exact integer geometry); 27\,195
networks and 5\,716 classes. Under 1~s.
\item \texttt{positive\_equilibria.py}: \zcref{prop:four}(2) for the 838 four-source nontrivial classes, i.e.
\(u,v,c,Q,N\), the 9 classes with \(N\equiv Q\equiv0\), the absence of sign change of \(QN\) on \(I\), the
zero structure of \(Q\), and the identification of the 30 classes with BBH's Table~1. About 10~s.
\item \texttt{classify\_all.py}: assigns every class one reason (trivial; at most two sources; three sources;
four sources with \(\Phi\) monotone or constant; Table~1 with inflow; Table~1 with an origin that is never AS,
checked symbolically in \(\kappa\); one of the ten). It also checks the totals and writes
\texttt{certificates/classes.csv}. About 10~s.
\item \texttt{certify\_ten.py}: the certificates of \zcref{tab:ten}, namely the equilibrium \((1,1)\), the signs of
\(\det\) and \(\Trop\), the origin, and the exact count of positive equilibria via \zcref{prop:four}(1) (roots of an
explicit polynomial isolated in rational intervals). About 10~s.
\item \texttt{explore\_numeric.py} (floating point, not part of the proof): for every nontrivial class with at least
three sources, it draws \(2\cdot10^4\) random rate constants with a prescribed positive equilibrium and tests for a
second AS equilibrium on the boundary. It finds hits in exactly the ten classes. About 16~minutes on 12 cores.
\item \texttt{index\_sanity.py} (floating point, not part of the proof): a check of \zcref{thm:index} on about
\(5\cdot10^4\) random quadratic fields \(F=Lz+q(z)\) with \(F_1(0,y)\ge0\), \(F_2(x,0)\ge0\) and a hyperbolic sink at the
origin. The sum of \(\operatorname{sign}\det DF\) over the positive equilibria is never positive. A control run with
a saddle at the origin produces positive sums, so the test is not vacuous. About 1~minute.
\end{itemize}

\paragraph{Limitations.} The proof relies on the exact computations of \zcref{prop:four}(2), which replicate BBH's
Theorem~25 and Remark~26(d), and on the classical facts quoted from~\autocite{KrasnoselskiiZabreiko1984}
and~\autocite{Perko2001}. The reading of the orange bullets of BBH's Table~1 was done on the rendered PDF.
It is consistent with the text of their Lemma~27 and with \texttt{classify\_all.py}: the ten classes with property
(MS) are exactly those marked. No verification has been carried out yet.

\paperreferences
\end{document}
